% Generated draft: requires mathematical and source review.
\documentclass[11pt]{article}
\usepackage[margin=25mm]{geometry}
\usepackage{fontspec}
\setmainfont{Latin Modern Roman}
\usepackage{amsmath,amssymb,mathtools,unicode-math}
\setmathfont{Latin Modern Math}
\usepackage{enumitem,needspace,xurl,hyperref}
\setlength{\parindent}{0pt}
\setlength{\parskip}{5pt}
\setlength{\emergencystretch}{3em}
\newcommand{\N}{\mathbb{N}}
\newcommand{\Z}{\mathbb{Z}}
\newcommand{\Q}{\mathbb{Q}}
\newcommand{\R}{\mathbb{R}}
\newcommand{\C}{\mathbb{C}}
\newcommand{\F}{\mathbb{F}}
\newcommand{\A}{\mathbb{A}}
\newcommand{\PP}{\mathbb{P}}
\newcommand{\E}{\mathbb{E}}
\newcommand{\eps}{\varepsilon}
\newcommand{\dd}{\,\mathrm{d}}
\begin{document}
\section*{UnsolvedMath ID: 8800012}
\textbf{Source problem code:} catalogue entry
\subsection{Definitions and mathematical statement}
For a variety \(X\) over a finite field \(\mathbb F_q\), its zeta function is the formal series
\[
 Z(X,T)=\exp\left(\sum_{r\ge1}\#X(\mathbb F_{q^r})\frac{T^r}{r}\right).
\]
In cohomological algorithms this is computed from Frobenius actions on suitable finite-dimensional cohomology groups. For a smooth proper setting where an appropriate characteristic-zero de Rham realization is available, cup product gives natural bilinear maps
\[
 H^i_{\mathrm{dR}}\times H^j_{\mathrm{dR}}\longrightarrow H^{i+j}_{\mathrm{dR}},
\]
and, after the top-degree trace, duality pairings. A \(p\)-adic implementation must specify its lift or comparison with the cohomology actually carrying Frobenius; arbitrary varieties are not assumed to admit a chosen smooth lift.

The source asks whether such natural pairings can improve zeta-function computations. A precise algorithmic target would identify the input class and encoding, baseline computation, coefficient/precision requirements, cost of constructing pairings, and the improvement sought: asymptotic time, memory, precision loss, certification, or practical efficiency. None is fixed in the excerpt, so no universal complexity bound or particular speedup is substituted. Using pairings for checks is distinct from proving an overall complexity improvement.

The supplied summary reports cup-product use in Kedlaya-style and certificate computations while saying a definitive improvement purely from cup products was not verified. It also includes unclear fragments ``google2024'', ``Klein/link'', and a ``Legendre symbol/certificate'' description, and names Roy--Saxena--Venkatesh (2024) without a precise publication. Those fragments do not define algorithms or establish the stated claims and need source checking. The 2006 workshop question and later partial-progress assessment remain attributions pending current primary-literature review.
\subsection{Short English statement}
Can natural de Rham cup-product pairings make finite-field zeta-function algorithms better? The variety class, cohomological model and intended improvement must be specified; the supplied summary's reported uses and modern references are incomplete.
\subsection{Sources}
\begin{itemize}
\item Local registry: \url{data/unsolvedmath/problems.json}, permanent UnsolvedMath ID \(8800012\), source code catalogue entry; supplied statement, background, and research summary, updated 2026-08-06. Attributed to Kiran Kedlaya.
\item IPAM, \emph{Workshop on Number Theory and Cryptography---Open Problems} (2006), attachment 6220, PDF pages 22--23, logical slide 10: \url{https://www.ipam.ucla.edu/programs/workshops/workshop-i-number-theory-and-cryptography-open-problems/?tab=schedule}. Supplied PDF/legacy-PPT visual audit accessed 2026-07-29; it reports inspecting 19 surviving attachments and excluding repetitions and solved definitions.
\item The summary mentions a Kedlaya survey using the unresolved label ``google2024'', a ``Klein/link'' fragment, and Roy--Saxena--Venkatesh (2024) without identifiable full citations or exact theorem statements.
\end{itemize}
\end{document}
